\chapter{Flow Dynamics of Delta–Sigma Loops}

The entropy-field interpretation introduced in the previous chapter reveals that the internal behavior of delta–sigma modulators is best described not as a discrete-time difference-equation system, but as a dynamical flow on a state manifold endowed with a vector field, curvature, and entropy density. In this chapter, we develop this view systematically. The result is a unified account of delta–sigma modulation as a flow-dynamic system whose stability, spectral behavior, and resolution arise from geometric and thermodynamic structure rather than from ad hoc frequency-domain arguments.

\section*{5.1 The State Manifold and its Intrinsic Geometry}

Let ( \mathcal{X} \subset \mathbb{R}^n ) denote the state space of an (n)-dimensional delta–sigma modulator. For a first-order modulator, ( n = 1 ); for higher-order architectures, ( n \ge 2 ). Although ( \mathcal{X} ) is typically regarded as a Euclidean vector space, the presence of quantization boundaries, loop-filter nonlinearities, saturation effects, and manifold constraints (introduced in Chapter~6) imply that the effective state space is curved and piecewise smooth. We therefore regard ( \mathcal{X} ) as a smooth manifold except at quantization boundaries, where it becomes stratified.

A state trajectory ( \gamma : \mathbb{Z} \rightarrow \mathcal{X} ) is generated by the iterative mapping
[
x[n+1] = F(x[n], u[n]),
]
where ( u[n] ) is the input. In the continuous-time limit, this becomes a flow generated by a vector field ( \mathbf{v} ):
[
\frac{dx}{dt} = \mathbf{v}(x(t),u(t)).
]
This vector field encodes the loop-filter dynamics and determines the geometry of the internal flow. The quantizer acts as a discontinuous operator that periodically resets the vector field by replacing ( u[n] ) with its discretized surrogate. The induced flow is therefore hybrid: smooth evolution is punctuated by discrete entropy-injection events.

Because the entropy field ( S(x,t) ) introduced previously evolves concurrently with the state, the pair ( (x(t), S(x,t)) ) evolves on the extended state manifold ( \mathcal{X} \times \mathbb{R}_{\ge 0} ), with the entropy component driven by injected uncertainty and transported by the vector field ( \mathbf{v} ).

\section*{5.2 The Vector Field Defined by Loop Filtering}

The loop filter defines the directional flow of the state. For concreteness, consider a general (p)-th order modulator with loop-filter matrix ( A \in \mathbb{R}^{p \times p} ) and input vector ( B \in \mathbb{R}^p ). The internal update equation is
[
x[n+1] = A,x[n] + B,(u[n] - y[n]).
]
In the continuous-time limit under standard bilinear or forward-Euler approximation, the associated vector field becomes
[
\mathbf{v}(x,u) = A,x + B,(u - Q(Cx)),
]
where (C) maps the state to the quantizer input.

This vector field defines a flow on ( \mathcal{X} ). The magnitude and direction of the flow at any point determine how quantization error is transported across the state manifold. For modulators with cascaded integrators (e.g., CIFF or CIFB topologies), the vector field contains nested directions corresponding to successive integrator states. The shape of the flow—its divergence, curl, and invariant sets—determines stability and shaping behavior. For example, if the divergence of the vector field is negative in a region, entropy tends to be compressed in that region; if positive, entropy spreads.

To formalize such properties, the divergence of ( \mathbf{v} ) is defined as
[
\nabla \cdot \mathbf{v}(x) = \sum_{i=1}^{p} \frac{\partial v_i}{\partial x_i}.
]
In classical linear cases with ( A ) stable and constant, the divergence is constant and negative, implying contraction. In nonlinear or adaptive modulators, the divergence varies across the state space, shaping entropy flows in complex ways.

\section*{5.3 Entropy Transport Under the Flow}

The entropy field ( S(x,t) ) evolves under the vector field ( \mathbf{v} ). The governing equation is the advection–diffusion equation with a source term:
[
\frac{\partial S}{\partial t} = -\nabla \cdot (S \mathbf{v}) + D \Delta S + \sigma(x,t),
]
where ( \sigma(x,t) ) represents the entropy injection due to quantization. The divergence term determines how entropy density is advected; the diffusion term represents smoothing induced by higher-order integration; the injection term introduces entropy when the quantizer collapses values into discrete cells.

In steady-state regime under constant input, the entropy field converges to a stationary solution satisfying
[
\nabla \cdot (S \mathbf{v}) = D \Delta S + \sigma(x).
]
This equation encapsulates delta–sigma noise shaping in a geometrically transparent way: the vector field transports entropy from low-frequency trajectories into higher-frequency ones, while diffusion (from higher-order filtering) smooths entropy gradients, preventing the system from developing catastrophic oscillations.

\section*{5.4 Stability as Flow Contraction}

Stability of a delta–sigma modulator corresponds to the non-escape of trajectories from the bounded region of the state manifold. In flow terms, the system is stable if all trajectories remain inside a compact set (K \subset \mathcal{X}). Let ( \phi_t(x_0) ) denote the flow generated by the vector field.

Stability requires that
[
\sup_{t \ge 0} | \phi_t(x_0) | < \infty.
]
A sufficient condition for this is that the divergence of the vector field be globally negative:
[
\nabla \cdot \mathbf{v}(x) < 0 \quad \forall x \in \mathcal{X}.
]
For linear cases, this requires that the loop-filter matrix (A) have eigenvalues strictly inside the unit circle (or left half-plane). For nonlinear or manifold-constrained systems, the divergence condition becomes local: stability may hold if the divergence is negative in the region containing all trajectories, even if positive elsewhere.

Lyapunov methods capture this through Lyapunov functions (V(x)) with negative directional derivatives:
[
\frac{d}{dt} V(x(t)) = \nabla V(x) \cdot \mathbf{v}(x) < 0.
]
In the entropy-field framework, a natural Lyapunov function includes both state energy and entropy:
[
V(x,S) = \frac{1}{2} |x|^2 + \alpha \int_{\mathcal{X}} S(x),dx.
]
A negative derivative indicates that both energy and entropy are being dissipated or advected away from the coherence manifold.

\section*{5.5 Limit Cycles as Flow Attractors}

Limit cycles arise when the vector field creates closed trajectories that trap the flow. In linear analysis, these correspond to periodic oscillations of the quantizer output. In the flow-dynamic view, they correspond to closed orbits on the state manifold where entropy is periodically injected and advected in a repeating pattern.

Let ( \Gamma \subset \mathcal{X} ) be a closed curve invariant under the flow. Then for initial conditions sufficiently near ( \Gamma ), the flow satisfies
[
\lim_{t\to\infty} \operatorname{dist}(\phi_t(x_0), \Gamma) = 0.
]
Such limit cycles correspond to frequency tones in the quantizer output and are often undesirable. They can be detected through Poincaré maps or through examining where the divergence changes sign, causing entropy accumulation at specific regions of the state manifold.

In higher-order modulators, such cycles may be distorted by curvature of the state manifold or disrupted by higher-dimensional flows that redistribute entropy more effectively. The presence of curvature (prefiguring Chapter~6) modifies the vector field through the connection coefficients and affects the stability of cycles.

\section*{5.6 Chaotic Regimes and Entropy Cascades}

Delta–sigma modulators can exhibit chaotic behavior, particularly when loop filters are aggressive or quantizer thresholds align pathologically with the vector flow. Chaos corresponds to a regime where small perturbations in the internal state lead to exponentially diverging trajectories.

The Lyapunov exponent ( \lambda ) satisfies
[
\lambda = \lim_{t\to\infty} \frac{1}{t} \log \frac{|\delta x(t)|}{|\delta x(0)|}.
]
Positive ( \lambda ) indicates chaos. In terms of the entropy field, chaos corresponds to the formation of entropy cascades wherein entropy is repeatedly injected and advected into increasingly fine scales, similar to turbulence in fluid dynamics. Indeed, the advection–diffusion structure of the entropy field admits an analogy to Kolmogorov cascades in turbulence theory \citep{frisch1995turbulence}.

Higher-order modulators are less susceptible to chaos because the diffusion operator ( D\Delta S ) smooths entropy gradients, preventing the formation of sharp entropy filaments. In Chapter~8, these smoothing mechanisms are analyzed using tools from stability theory and modern control.

\section*{5.7 The Flow Interpretation of Noise Shaping}

Classical frequency-domain analysis asserts that the noise transfer function ( \mathrm{NTF}(z) ) suppresses low-frequency quantization noise by forcing the error signal through a high-pass filter. In the flow-dynamic formulation, this effect is understood geometrically. The vector field ( \mathbf{v} ) directs entropy away from the low-frequency stable manifold and towards unstable or oscillatory subspaces. These subspaces correspond to high-frequency modes.

Let ( \mathcal{M}*{\text{coh}} \subset \mathcal{X} ) denote the coherence manifold, the subspace where low-frequency structure is preserved. Let ( \mathcal{M}*{\text{ent}} \subset \mathcal{X} ) denote the entropy manifold, the subspace where uncertainty accumulates. The flow decomposes ( \mathcal{X} = \mathcal{M}*{\text{coh}} \oplus \mathcal{M}*{\text{ent}} ). Projection of the quantization error (e[n]) into these subspaces yields
[
e[n] = P_{\text{coh}} e[n] + P_{\text{ent}} e[n].
]
The delta–sigma loop is designed such that
[
P_{\text{coh}} e[n] \approx 0,
]
and
[
P_{\text{ent}} e[n] \approx e[n],
]
meaning low-frequency components of the error vanish while high-frequency ones dominate. Frequency-domain noise shaping emerges as the spectral footprint of the entropy-transport process.

\section*{5.8 The Continuous-Time Limit and PDE Correspondence}

The continuous-time limit of a delta–sigma modulator with quantization noise treated as a rapidly varying perturbation yields a system of coupled PDEs governing the joint evolution of the signal and entropy fields:
[
\frac{dx}{dt} = \mathbf{v}(x, u, S),
]
[
\frac{\partial S}{\partial t} = -\nabla \cdot (S \mathbf{v}) + D \Delta S + \sigma(x,t).
]
This PDE structure makes the analogy with RSVP explicit. In RSVP, scalar fields ( \Phi(x,t) ), vector fields ( \mathbf{v}(x,t) ), and entropy fields ( S(x,t) ) coevolve via transport, smoothing, and constraint equations. Delta–sigma modulation represents a discrete symbolic shadow of this continuous flow.

\section*{5.9 Transition to Manifold-Aligned Representations}

The vector-field formulation developed in this chapter depends critically on the geometry of the state manifold ( \mathcal{X} ). If the signal lies on a low-dimensional manifold embedded in a higher-dimensional space, the vector field should be aligned with that manifold to minimize error and avoid artificial instability. The next chapter develops this idea rigorously by introducing tangent spaces, geometric projection, and curvature effects, showing how delta–sigma modulators gain stability and accuracy by respecting the intrinsic geometry of the signals they process.

